\chapter*{\huge Remerciements}
\markboth{Remerciements}{Remerciements}
\vspace{6mm}

Ce mémoire porte une signature, mais il doit son existence à un ensemble de personnes qui, chacune
à sa place, ont rendu ce travail possible. Qu'elles trouvent ici l'expression de ma reconnaissance.

Je tiens à remercier en premier lieu \textbf{Monsieur Maher DAASSI}, Directeur des Systèmes
d'Information d'OneTech Business Solutions, pour la confiance qu'il m'a accordée en ouvrant les
portes de son service à un projet dont personne ne pouvait garantir l'issue au moment où il a été
lancé. Son soutien ne s'est jamais démenti, et le temps qu'il a bien voulu me consacrer, malgré des
responsabilités qui laissent peu de place à ce genre de disponibilité, a compté davantage qu'il ne
l'imagine. Sa vision du métier m'a permis de comprendre ce qu'un service informatique attend
réellement d'un outil, et de distinguer ce qui relève de la prouesse technique de ce qui apporte
une valeur mesurable au quotidien. Plusieurs orientations décisives de ce travail lui reviennent
directement.

J'exprime toute ma gratitude à \textbf{Madame Rahma OUESLATI}, Consultante en Systèmes
d'Information, mon encadrante professionnelle. Elle m'a accompagné du premier jour jusqu'à la
livraison, avec une disponibilité constante et une exigence bienveillante. Ses conseils techniques
m'ont épargné plus d'une impasse, et sa connaissance fine du traitement réel des tickets a nourri
chaque décision de conception. Elle a su me laisser chercher, tout en me remettant sur la voie
lorsque je m'en éloignais.

Je remercie également \textbf{Monsieur Ahmed DHOUIBI}, mon encadrant académique, pour la rigueur
de son suivi et la justesse de ses remarques. Son regard critique m'a conduit à reprendre certains
choix que je croyais acquis, et à assumer explicitement les résultats qui n'ont pas répondu aux
objectifs initiaux plutôt qu'à les contourner. C'est probablement ce que ce mémoire doit de plus
précieux à son encadrement.

Mes remerciements s'adressent aussi à \textbf{l'ensemble des membres et collaborateurs d'OneTech
Business Solutions}, et particulièrement aux techniciens du service support, pour la qualité de
leur accueil. Ils ont accepté qu'un stagiaire observe leur travail, pose des questions parfois
naïves, et propose de l'outiller. Leur ouverture et le cadre de travail qu'ils m'ont offert ont
fait de ces mois une expérience professionnelle autant qu'un projet de fin d'études.

Je tiens enfin à remercier \textbf{le corps enseignant et l'équipe pédagogique de TEK-UP
University} pour la formation reçue durant ces années, ainsi que \textbf{les membres du jury} qui
me font l'honneur d'évaluer ce travail.
